\chapter{Análisis de Resultados}
\label{cha:analisis_resultados}

En esta sección se presentan los resultados obtenidos a partir de la implementación y ejecución del modelo Transformer basado en roBERTa explicado en \ref{sec:RoBERTa} para el etiquetado de sugerencias sobre el dataset de reseñas turísticas de \textcite{tripadvisor-website}. Realizaremos un análisis mixto (cualitativo y cuantitativo) de los resultados obtenidos, evaluando datos cuantitativos como la precisión del modelo y datos cualitativos como calidad de las sugerencias identificadas.

Habiendo realizado una larga investigación y luego experimentaciones con distintos modelos NLU, realizando hipótesis de su posible rendimiento teóricos para la tarea y analizando su comportamiento en la práctica demostrando de manera empírica su funcionamiento real, podemos llegar a ciertas conclusiones respecto a estos mismos, a como fue realizado el trabajo en general y mas importante, hasta donde llegó y hacia donde se puede dirigir esta investigación (aunque eso sea detallado mas a fondo en la sección de trabajo futuro \ref{cha:trabajo_futuro}).

Podemos concluir que el modelo RoBERTa Base después del proceso de fine-tuning ha logró obtener un buen rendimiento en la clasificación de los comentarios de turistas, acercandose en un 80\% en precisión a la clasificación que fue brindada por el modelo OLMO 3 THINKING. Además, el modelo logró obtener un buen equilibrio entre precisión y recall, lo que se refleja en un F1-score de 0.774, lo que indica que el modelo fue capaz de identificar correctamente tanto las sugerencias como las no sugerencias con una tasa de acierto aceptable. Sin embargo, es importante mencionar que el modelo aún tiene margen de mejora, especialmente en la identificación de sugerencias, ya que el recall es de 0.748, lo que indica que hay un porcentaje de sugerencias no menor que el modelo no está identificando correctamente. Esto podría ser debido a la cantidad limitada de datos de sugerencias disponibles para el entrenamiento, como también, a los fallos que el modelo OLMO 3 THINKING pudo haber cometido en el proceso de etiquetado, lo que podría haber introducido ruido en el dataset de nivel de plata utilizado para el entrenamiento del modelo RoBERTa Base. 

A pesar de esto hemos podido presenciar una mejora significativa en los tiempos de ejecución del modelo RoBERTa Base en comparación con el modelo OLMO 3 THINKING, ya que el modelo RoBERTa al ser mucho más pequeño y no requiere tokens extras para el razonamiento y análisis de los comentarios, hemos logrado reducir el tiempo de ejecución etiquetando 10000 comentarios en 118 horas con OLMO 3 THINKING a tan solo 30 minutos con el modelo RoBERTa Base ejecutando el programa en un Macbook Air con chip M4, 16 GB de RAM, 10 Nucleos de CPU y 10 Nucleos de GPU, lo que representa una mejora significativa en términos de eficiencia y escalabilidad para el proceso de clasificación de los comentarios turistas.

% =================================================================
% = Tabla comparativa de modelos – métricas en conjunto de prueba =
% =================================================================

Podemos observar gracias a la tabla \ref{fig:Comparativa de modelos frente a RoBERTa Base en el conjunto de prueba} que el modelo RoBERTa Base logró superar a todos los modelos de manera satisfactoria en terminos de las métricas utilizadas para la comparación de los mismos, existiendo la mayor diferencia en los resultados de la precisión con una diferencia de 0.018 (1.8\%) con respecto al siguiente en segundo lugar.
